% =========================================================
% Volume IV  \S{}1.7.5  --- Q as a Model of the Field Signature
% WIRING NOTE (cross-volume):
%   The \hyperref targets thm:q-is-field, thm:q-is-integral-domain,
%   thm:q-is-ordered-field, def:algebraic-field,
%   def:addition-on-q-classes, def:multiplication-on-q-classes,
%   def:positive-cone live in Volume II (rationals chapter) and resolve
%   in the assembled monorepo build (Learning-Real-Analysis/main.tex).
%   thm:z-embeds-in-q, thm:q-embeds-in-r live in Volume II Ch.8
%   (Embedding the Number Systems) once that chapter is merged.
%   In a STANDALONE volume-iv build these resolve to '??'; that is the
%   expected, honest behavior of a cross-volume reference, not a bug.
% =========================================================
\subsection*{The Rational Numbers as a Structure}

\begin{remark*}[Modeling Goal]
We exhibit $\mathbb{Q}$ as a concrete \emph{model}: fix the field signature,
take the carrier and operations built in Volume~II, interpret each symbol, and
discharge satisfaction by citing the Volume~II theorem that $\mathbb{Q}$ obeys
the field axioms. Order is then added by \emph{expansion}, and the additive
group is recovered by \emph{reduction}.
\end{remark*}

% ---- 1. the signature being interpreted ----
\begin{remark*}[Exposition]
A model is a signature together with an interpretation on a carrier. The
signature here is the field fragment of the
\hyperref[def:arithmetic-signature]{arithmetic signature},
\[
\mathcal{L}_{\mathrm{field}}=\{+,\cdot,-,{}^{-1},0,1\},
\]
with binary function symbols $+,\cdot$, a unary $-$, a unary ${}^{-1}$ defined on
nonzero elements, and constant symbols $0,1$. The order symbol $<$ is
deliberately absent; it enters by expansion below.
\end{remark*}

% ---- 2. the model itself ----
\begin{definition}[The Rational Field Model $\mathcal{Q}$]
\label{def:rational-field-model}
The \textbf{rational field model} is the $\mathcal{L}_{\mathrm{field}}$-structure
\[
\mathcal{Q}=\bigl(\mathbb{Q};\,
+^{\mathcal{Q}},\,\cdot^{\mathcal{Q}},\,-^{\mathcal{Q}},\,{}^{-1\,\mathcal{Q}},\,
0^{\mathcal{Q}},\,1^{\mathcal{Q}}\bigr),
\]
whose carrier $|\mathcal{Q}|=\mathbb{Q}$ is the set of rational numbers
constructed in Volume~II, and whose symbols are interpreted by the
\hyperref[def:addition-on-q-classes]{rational addition} and
\hyperref[def:multiplication-on-q-classes]{rational multiplication}: $+^{\mathcal{Q}}$ and
$\cdot^{\mathcal{Q}}$ are rational addition and multiplication, $-^{\mathcal{Q}}$
is rational negation, ${}^{-1\,\mathcal{Q}}$ is the inverse of a nonzero
rational, and $0^{\mathcal{Q}},1^{\mathcal{Q}}$ are the rational constants
$0$ and $1$.
\end{definition}

\begin{remark*}[Standard quantified statement]
\[
|\mathcal{Q}|=\mathbb{Q},\qquad
+^{\mathcal{Q}},\cdot^{\mathcal{Q}}:\mathbb{Q}\times\mathbb{Q}\to\mathbb{Q},\qquad
-^{\mathcal{Q}}:\mathbb{Q}\to\mathbb{Q},\qquad
0^{\mathcal{Q}},1^{\mathcal{Q}}\in\mathbb{Q}.
\]
\end{remark*}\begin{remark*}[Interpretation]
Nothing here is new mathematics: the carrier and operations are exactly the ones
built in Volume~II. The model only packages them as a single object
$\mathcal{Q}$, so that ``$\mathbb{Q}$'' can be handed to any theorem that
quantifies over fields, and so the satisfaction statement below has a concrete
structure to range over. The signature is the type; $\mathcal{Q}$ is the
instance.
\end{remark*}

% ---- 3. satisfaction: the keystone, wired to Volume II ----
\begin{remark*}[Satisfaction]
$\mathcal{Q}$ is a \emph{model} of the field axioms:
\[
\mathcal{Q}\models \Sigma_{\mathrm{field}},
\qquad\text{equivalently}\qquad
\operatorname{Field}(\mathbb{Q},+,\cdot,0,1).
\]
This is not re-proved here. It is exactly the content of
\hyperref[thm:q-is-field]{$\mathbb{Q}$ Is a Field} in Volume~II, which verifies
every field axiom against the rational operations; that theorem is therefore the
witness that the interpretation $\mathcal{Q}$ satisfies
$\Sigma_{\mathrm{field}}$. The integral-domain fragment is witnessed separately
by \hyperref[thm:q-is-integral-domain]{$\mathbb{Q}$ Is an Integral Domain}.
\end{remark*}

% ---- 4. expansion to the ordered field ----
\begin{remark*}[Expansion to the ordered field]
Adjoining the relation symbol $<$ gives the expansion
\[
\mathcal{Q}_{<}=(\mathbb{Q};+,\cdot,-,{}^{-1},0,1,<),
\qquad
\operatorname{Expansion}(\mathcal{Q}_{<},\mathcal{Q}),
\]
in the sense of \hyperref[def:structure-expansion]{Structure Expansion}: the
carrier is unchanged and exactly one new interpreted symbol is added, where
$<^{\mathcal{Q}}$ is the rational order determined by the
\hyperref[def:positive-cone]{positive cone}. That $\mathcal{Q}_{<}$ satisfies the
ordered-field axioms is witnessed by
\hyperref[thm:q-is-ordered-field]{$\mathbb{Q}$ Is an Ordered Field} in
Volume~II.
\end{remark*}

% ---- 5. reduct to the additive group ----
\begin{remark*}[Reduct to the additive group]
Forgetting $\cdot,{}^{-1},1$ leaves the reduct
\[
\mathcal{Q}\!\restriction\!\{+,-,0\}=(\mathbb{Q};+,-,0),
\qquad
\operatorname{Reduction}\!\bigl(\mathcal{Q}\!\restriction\!\{+,-,0\},\,\mathcal{Q}\bigr),
\]
in the sense of \hyperref[def:structure-reduction]{Structure Reduction}: the same
carrier with fewer visible symbols. This reduct is the additive abelian group of
$\mathbb{Q}$ --- what remains when only additive structure is in view.
\end{remark*}

% ---- 6. blueprint position + embedding chain (wire to Vol II Ch.8) ----
\begin{remark*}[Blueprint position]
In the structural blueprint $\mathcal{Q}$ is the \emph{field of fractions} of the
integer ring: the smallest field into which the integer model embeds. This
places it on the chain
\[
\mathbb{Z}\hookrightarrow\mathbb{Q}\hookrightarrow\mathbb{R},
\]
whose arrows are the canonical embeddings of Volume~II
(\hyperref[thm:z-embeds-in-q]{$\mathbb{Z}$ Embeds in $\mathbb{Q}$},
\hyperref[thm:q-embeds-in-r]{$\mathbb{Q}$ Embeds in $\mathbb{R}$}).
Model-theoretically the first embedding presents the integer model as a
substructure of the $\{+,\cdot,0,1\}$-reduct of $\mathcal{Q}$; passing to the
full field $\mathcal{Q}$ supplies the multiplicative inverses the integer model
lacks.
\end{remark*}

\begin{remark*}[Interpretation]
$\mathbb{Q}$ is the first number system whose model satisfies the \emph{full}
field signature: $\mathbb{Z}$ is only a ring (no inverses), and $\mathbb{Q}$
closes exactly that gap. Adding $<$ by expansion gives the ordered field; adding
analytic completeness --- which is \emph{not} a first-order property of the field
signature --- is deferred to the real model in the analysis layer.
\end{remark*}

\begin{dependencies}
\begin{itemize}
  \item \hyperref[def:arithmetic-signature]{Arithmetic Signature}.
  \item \hyperref[def:algebraic-field]{Field} (Volume~II).
  \item \hyperref[def:rationals]{The Rational Numbers}.
  \item \hyperref[def:addition-on-q-classes]{Rational Addition}.
  \item \hyperref[def:multiplication-on-q-classes]{Rational Multiplication}.
\end{itemize}
\end{dependencies}
